% ============================================================================
% 问题三 · 建模与求解
%
% 【四问通用合同】本片是全文体量最大的部分，必须写透这五段：
%   1. 模型建立：变量、目标、约束、方程。推导只留关键结论，全过程移附录。
%   2. 求解算法：复杂度、分辨率预算、样本量及其来源（怎么定的，不是拍的）。
%   3. 主结果：数字全部走宏，禁止手抄。
%   4. 证书：置信下界 / 收敛判据 / 双路互证，至少一样。只给点估计的结论要降措辞。
%      最优性主张分级：global / local / best-found 三选一，无据不得写「最优」。
%   5. 图表解读：每张图配一段「从图上看出什么」，不是只把图放上去。
%
% 【搜索类问题的额外红线】网格步长必须由可行前沿的陡度决定：
%   前沿附近退化为单位步长，或把由别问解出的阈值点强制塞进候选集。
%   两家候选都栽在「网格把最优点跨过去」上，这是高频失分点。
% 篇幅参考：每问 3-6 页。
% ============================================================================
\subsubsection{模型的建立}

\subsubsection{求解算法与实现}

\subsubsection{结果与统计证书}

\subsubsection{结果讨论}
